\section{AI智能体在资料归集场景中的应用实践}

\subsection{典型场景一：站点证照与影像资料的自动归集}

在实际工作中，笔者曾处理过一项具有代表性的资料整理任务。该任务面向新源片区共计 19 个站点，原始目录中包含两类文件夹和大量散落照片。其中，两类文件夹分别为不动产证资料和营业执照正本资料；其余则为 19 个站点对应的便利店门头照片、便利店内景照片以及《仅销售预包装食品经营者备案信息采集表》照片。除前述两个目录外，当前目录下其余文件均为 JPG 格式照片，共计 57 张。该批资料在来源、文件命名和存放结构上均存在明显差异，尤其是部分照片命名极不统一，难以通过传统文件检索方式直接完成批量归档。

类似的问题如果完全依靠人工处理，工作人员通常需要反复打开图片或文档，识别分类依据信息，再按照目标结构逐项建立文件夹、复制剪切文件并进行校验。这一过程不仅机械，而且极易因命名相近、文件数量多而发生遗漏或误放。若引入 AI 智能体，则可通过“站点识别+资料类别判断+目标路径映射”的规则体系，实现批量化整理。其结果不是简单地“做了这件事”，而是通过智能体的介入，形成了一个可复用的业务能力模块，能够在未来类似场景中快速复制和推广。

\subsubsection{整理前后的目录结构示意}

为更直观地呈现 AI 智能体介入前后的差异，本文以目录树形式展示原始资料结构与目标归档结构。整理前的特点是“少量目录+大量散落照片并存”；整理后的特点则是“按站点统一归集，同一站点下集中放置相关资料”。

\newpage

\begin{lstlisting}[style=treeview,caption={整理前的示意性目录结构}]
资料原始目录/
├── 不动产证/
│   ├── 卡普河站不动产/
│   ├── 天河站不动产证/
│   └── ...
├── 营业执照正本/
│   ├── 营业执照正本~卡普河站.jpg
│   ├── 天河站营业执照正本.jpg
│   └── ...
├── 便利店门头~卡普河加油站.jpg
├── 便利店内~卡普河加油站.jpg
├── 天河站便利店门头.jpg
├── 天河站便利店内.jpg
├── 仅销售预包装食品经营者备案信息采集表~卡普河站.jpg
├── 天河站仅销售预包装食品经营者备案信息采集表.jpg
└── ...
\end{lstlisting}

\begin{lstlisting}[style=treeview,caption={整理后的示意性目录结构}]
站点资料标准目录/
├── 卡普河站/
│   ├── 不动产证_卡普河站/
│   ├── 营业执照正本_卡普河站.jpg
│   ├── 便利店门头_卡普河站.jpg
│   ├── 便利店内_卡普河站.jpg
│   └── 仅销售预包装食品经营者备案信息采集表_卡普河站.jpg
├── 天河站/
│   ├── 不动产证_天河站/
│   ├── 营业执照正本_天河站.jpg
│   ├── 便利店门头_天河站.jpg
│   ├── 便利店内_天河站.jpg
│   └── 仅销售预包装食品经营者备案信息采集表_天河站.jpg
└── ...
\end{lstlisting}

\subsection{典型场景二：站点考勤表的批量识别、重命名与缺报核验}

除静态资料归集外，AI 智能体在月度经营管理资料的批量整理场景中同样具有较强的适用性。以人员考勤资料为例，站点上报文件往往存在命名不统一、站点简称混用、合同化与劳务派遣材料并存等问题，导致后续汇总、追踪与催报工作需要投入额外精力。针对这一类场景，智能体不仅需要完成重命名，还需要结合站点台账判断哪些站点尚未提交材料，并输出便于后续追踪的缺报清单。

本文选取新源片区 \verb|D:/工作/CNPC/片区/新源片区/综合岗/综合办公室/人员考勤/202603/| 目录作为演示路径，采用命令行方式调用智能体执行任务。演示要求包括：一是将当前目录下的 Excel 文件统一重命名为 \lstinline|"202603_xxx加油站考勤统计表"| 或 \lstinline|"202603_xxx加气站考勤统计表"| 的格式；二是在同一站点同时存在合同化员工与劳务派遣员工考勤表的情况下，确保两个文件不会被错误覆盖，并在必要时通过表格结构或表内内容进一步判断；三是对于 \lstinline|"考勤表（2025版）.xlsx"| 这类文件名无法直接体现站点信息的材料，允许智能体打开 Excel，依据表头、站点名称栏或签章信息识别其实际对应站点；四是读取 \lstinline|"站点+站长.txt"| 中的站点台账，在任务执行完毕后自动生成或覆写一份 \lstinline|"未提交站点.txt"|；五是沉淀一份可复用的 \lstinline|AGENTS.md| 操作规范，并在完成后删除原始 \lstinline|"站点+站长.txt"| 文件。

\begin{codeexample}[htbp]
   \caption{命令行演示与任务指令}
   \begin{lstlisting}[style=cmdview]
D:\工作\CNPC\片区\新源片区\综合岗\综合办公室\人员考勤\202603>codex --yolo

重命名当前目录下的Excel文件，格式为：
"202603_xxx加油站考勤统计表"或"202603_xxx加气站考勤统计表"

要求：
1. xxx在原文件中会有体现
2. 区分同一站点合同化员工和劳务派遣员工，不能导致两个相异的文件被重复命名；
   如若无法区分，可以查看Excel表结构和内容进行判断。
3. "站点+站长.txt"是新源片区所有站点和站长的文件，执行完参考这个文件生成
   或覆写一份"未提交站点.txt"文件以供我查询谁未提交
4. 执行完所有指令之后，生成一份AGENTS.md文件，固定规范复用执行；
   "站点+站长.txt"的内容同样写入md文件；最后删除"站点+站长.txt"
\end{lstlisting}
\end{codeexample}

\subsubsection{整理前后的目录结构示意}

在该场景下，原始目录并非站点维度的结构化资料，而是多个文件名风格并存的散乱目录。通过智能体执行后，文件将按照统一规则完成识别与重命名，并额外生成 \lstinline|"未提交站点.txt"| 用于缺报追踪。为便于对比，本文采用上下顺排方式展示整理前后的目录状态。

\begin{lstlisting}[style=treeview,caption={整理前的示意性目录结构}]
伊犁公司新源县加气站2026年3月合同化员工考勤表.xlsx
伊犁公司阿勒玛勒加油站2026年3月合同化员工考勤表（2025年新版).xlsx
副本铁新公路站-新源片区2026年3月考勤统计表（202408新版）.xlsx
新源县加气站2026年3月劳务派遣考勤表.xlsx
新源片区阿勒玛勒2026年3月考勤统计表（202408新版）.xlsx
木斯服务区北加油站站2026年3月考勤统计表（202408新版）.xlsx
木斯服务区南加油站站2026年3月考勤统计表（202408新版）.xlsx
种羊场加气站员工考勤表-2026年3月.xlsx
种羊场加油站2026年3月考勤统计表（202408新版）.xlsx
考勤表（2025版）.xlsx
那拉提高速北加油站2026年3月考勤统计表.xlsx
那拉提高速南加油站2026年3月考勤统计表-.xlsx
站点+站长.txt
\end{lstlisting}

\newpage

\begin{lstlisting}[style=treeview,caption={整理后的示意性目录结构}]
202603_卡普河站考勤统计表.xlsx
202603_阿勒玛勒站合同化员工考勤统计表.xlsx
202603_阿勒玛勒站劳务派遣员工考勤统计表.xlsx
202603_铁新加油站考勤统计表.xlsx
202603_新源县加气站合同化员工考勤统计表.xlsx
202603_新源县加气站劳务派遣员工考勤统计表.xlsx
202603_木斯服务区北站考勤统计表.xlsx
202603_木斯服务区南站考勤统计表.xlsx
202603_种羊场加气站考勤统计表.xlsx
202603_种羊场站考勤统计表.xlsx
202603_那拉提服务区北站考勤统计表.xlsx
202603_那拉提服务区南站考勤统计表.xlsx
未提交站点.txt
AGENTS.md
\end{lstlisting}

\subsubsection{执行结果与缺报核验}

从执行结果看，智能体在该场景中完成了三项关键任务。其一，对文件名中 \lstinline|"新源片区"|、\lstinline|"副本"|、\lstinline|"202408新版"|、\lstinline|"2025版"| 等非核心表述进行了剥离，保留了月份、站点和资料主题等关键字段，使文件命名回归统一规则。其中，原始文件 \lstinline|"考勤表（2025版）.xlsx"| 并未在文件名中直接给出站点名称，智能体需要实际读取 Excel 内容，并根据表内出现的“卡普河加油站”等关键信息判断其对应站点，最终将其规范化为 \lstinline|"202603_卡普河站考勤统计表.xlsx"|。其二，对同一站点的合同化员工考勤表与劳务派遣员工考勤表进行了区分处理，避免出现因同名覆盖导致的数据丢失。其三，结合 \lstinline|"站点+站长.txt"| 中的完整站点台账，生成了 \lstinline|"未提交站点.txt"|，将缺报情况从人工逐一核对转化为自动比对输出。

按照新源片区站点台账，共有 19 个站点纳入比对范围：卡普河站、阿热勒托别站、城北站、则克台镇西站、塔勒德镇西站、则克台镇东站、阿勒玛勒站、哈拉布拉站、种羊场站、天河站、木斯站、新河站、木斯服务区南站、木斯服务区北站、那拉提服务区南站、那拉提服务区北站、铁新加油站、新源县加气站、种羊场加气站。结合演示目录中的已提交材料，可以判定本次未提交站点包括：阿热勒托别站、城北站、则克台镇西站、塔勒德镇西站、则克台镇东站、哈拉布拉站、天河站、新河站、木斯站。由此可见，AI 智能体并非只是在“把文件名改整齐”，而是能够进一步完成站点映射、材料去重、异常识别和缺报输出等成套动作，体现出较强的执行稳定性与业务可靠性。

\subsection{AI智能体驱动的业务价值}

从业务效果看，该类应用至少带来三方面收益。其一，显著降低资料归集所需的重复劳动时间，使工作人员从机械性整理中释放出来。其二，通过统一命名与目录结构，提升资料管理的一致性和可检索性，减少因个人习惯差异造成的后续沟通成本。其三，为经营资料进一步进入结构化管理、知识沉淀和智能分析创造前提条件。换言之，AI 智能体在此类场景中的意义，不仅在于“替人做事”，更在于把原本松散的文件处理过程，转化为可复用、可复制、可扩展的标准化业务能力。

因此，在基层经营管理实践中，AI 智能体并不是附加性的技术装饰，而是推动经营资料管理从经验驱动走向规则驱动、从人工分散走向流程集成的重要抓手。围绕资料归集、数据清洗、报表分析和文本生成等环节持续推进智能化建设，有望形成可推广的数智化办公范式。
